\section*{Chapter \ref{ch:g12:space} --- Vectors, Lines and Planes in Space}

\begin{solution}{exo:g12:space:1}
$\vect{AB}(2, 1, -1)$ and $\vect{AC}(1, -1, 1)$ are not proportional
($\frac21 \neq \frac{1}{-1}$), so the points are not aligned. Line $(AB)$:
\[
(x, y, z) = (1 + 2t,\; t,\; 2 - t), \qquad t \in \R .
\]
\end{solution}

\begin{solution}{exo:g12:space:2}
$2(x - 1) - (y - 1) + 3(z - 0) = 0$, \emph{i.e.}
\[
2x - y + 3z - 1 = 0 .
\]
For $B(0,2,1)$: $0 - 2 + 3 - 1 = 0$: yes, $B$ lies on the plane.
\end{solution}

\begin{solution}{exo:g12:space:3}
$\vect{AB}(1,1,0)$, $\vect{AC}(1,0,1)$:
\[
\cos\widehat{A}
= \frac{\vect{AB}\cdot\vect{AC}}{\norm{\vect{AB}}\,\norm{\vect{AC}}}
= \frac{1}{\sqrt2\,\sqrt2} = \frac12,
\]
so $\widehat A = 60^\circ$.
\end{solution}

\begin{solution}{exo:g12:space:4}
Points of $\mathcal D$: $(1 + t,\, 2 - t,\, 2t)$. Substituting into the
\omterm{def:g10:algebra:equation}{equation} of $\mathcal P$:
\[
2(1+t) + (2 - t) - 2t + 1 = 5 - t = 0 \iff t = 5 .
\]
Unique \omterm{def:g10:numbers:interunion}{intersection} point: $(6, -3, 10)$.
\end{solution}

\begin{solution}{exo:g12:space:5}
The \omterm{def:g12:space:normal}{normal vectors} $(1, -1, 2)$ and $(2, 1, 1)$ are not \omterm{def:g12:space:collinear}{collinear}, so the
planes intersect in a line. Adding the two \omterm{def:g10:algebra:equation}{equations}: $3x + 3z = 5$, so
$x = \frac53 - z$; then the first \omterm{def:g10:algebra:equation}{equation} gives
$y = x + 2z - 1 = \frac23 + z$. With $z = t$:
\[
(x, y, z) = \left(\frac53 - t,\; \frac23 + t,\; t\right), \qquad t \in \R,
\]
a line with \omterm{def:g11:vect:direction}{direction vector} $(-1, 1, 1)$. (Both \omterm{def:g10:algebra:equation}{equations} check
identically.)
\end{solution}

\begin{solution}{exo:g12:space:6}
Directions $\vec u_1(1, 2, -1)$ and $\vec u_2(1, 1, 2)$ are not \omterm{def:g12:space:collinear}{collinear},
so the lines are not parallel. Intersecting would require
\[
1 + t = 2 + s, \qquad 2t = 1 + s, \qquad 3 - t = 1 + 2s .
\]
The first two give $t = 1 + s$ and $2(1+s) = 1 + s$, so $s = -1$, $t = 0$;
but then the third reads $3 = -1$, absurd. No common point and not
parallel: the lines are skew.
\end{solution}

\begin{solution}{exo:g12:space:7}
\emph{1.} $\operatorname{dist} = \dfrac{\abs{2\cdot3 - 1 + 2\cdot1 - 3}}
{\sqrt{4 + 1 + 4}} = \dfrac{4}{3}$.

\emph{2.} $H = M_0 + t\,\vec n$ with $\vec n(2, -1, 2)$, choosing $t$ so
that $H \in \mathcal P$:
\[
2(3 + 2t) - (1 - t) + 2(1 + 2t) - 3 = 4 + 9t = 0 \iff t = -\frac49 .
\]
Hence $H\left(\frac{19}{9}, \frac{13}{9}, \frac19\right)$, and
\[
M_0H = \norm{t\,\vec n} = \frac49 \times 3 = \frac43,
\]
matching the distance formula.
\end{solution}

\begin{solution}{exo:g12:space:8}
\omterm{def:g10:coordgeom:system}{Coordinates}: $G(1,1,1)$, and the plane $(BDE)$ passes through $B(1,0,0)$,
$D(0,1,0)$, $E(0,0,1)$.

\emph{1.} The plane $(BDE)$ has \omterm{def:g10:algebra:equation}{equation} $x + y + z = 1$ (satisfied by the
three points, and they are not aligned), so $\vec n(1,1,1)$ is normal to
it. Since $\vect{AG} = (1,1,1) = \vec n$, the diagonal $(AG)$ is \omterm{def:g11:scal:orthogonal}{orthogonal}
to the plane $(BDE)$.

\emph{2.} Distance from $A(0,0,0)$:
$\frac{\abs{0 + 0 + 0 - 1}}{\sqrt3} = \frac{1}{\sqrt3} = \frac{\sqrt3}{3}$.
The line $(AG)$ is $(t, t, t)$; it meets the plane when $3t = 1$: at the
point $\left(\frac13, \frac13, \frac13\right)$, the centroid of the
triangle $BDE$.
\end{solution}

\begin{solution}{exo:g12:space:9}
\emph{1.} $\vect{AB}(1, 0, -1)$, $\vect{AC}(-1, 1, 0)$,
$\vect{AD}(1, 1, 1)$. If $\vect{AD} = a\vect{AB} + b\vect{AC}$, the
\omterm{def:g10:coordgeom:system}{coordinates} give $a - b = 1$, $b = 1$, $-a = 1$: the first two force
$a = 2$, contradicting $a = -1$. Not \omterm{def:g12:space:collinear}{coplanar}.

\emph{2.} $\vect{BC}(-2, 1, 1)$, $\vect{BD}(0, 1, 2)$. Solving
$\vec n \cdot \vect{BC} = \vec n \cdot \vect{BD} = 0$: from
$b + 2c = 0$, $b = -2c$; then $-2a - 2c + c = 0$ gives $c = -2a$. With
$a = 1$: $\vec n(1, 4, -2)$. Plane through $B(2,1,0)$:
\[
x + 4y - 2z - 6 = 0 .
\]

\emph{3.} Distance from $A(1,1,1)$:
$\dfrac{\abs{1 + 4 - 2 - 6}}{\sqrt{21}} = \dfrac{3}{\sqrt{21}}$.
Area of $BCD$: $\norm{\vect{BC}}^2 = 6$, $\norm{\vect{BD}}^2 = 5$,
$\vect{BC}\cdot\vect{BD} = 3$, so
\[
\text{Area} = \frac12\sqrt{6 \times 5 - 9} = \frac{\sqrt{21}}{2}.
\]

\emph{4.}
\[
V = \frac13 \times \frac{\sqrt{21}}{2} \times \frac{3}{\sqrt{21}}
= \frac12 .
\]
\end{solution}

\begin{solution}{pb:g12:space:1}
\textbf{1.} Plane: $x + y + z = 1$. Line: $(t, t, t)$.
\omterm{def:g10:numbers:interunion}{Intersection}: $3t = 1$:
$\left(\frac13, \frac13, \frac13\right)$.

\textbf{2.} $\dfrac{\abs{0 + 0 + 0 - 1}}{\sqrt3} =
\dfrac{1}{\sqrt3} = \dfrac{\sqrt3}{3}$.

\textbf{3.} $1 - 1 + 0 = 0$: \omterm{def:g11:scal:orthogonal}{orthogonal}.

\textbf{4.} Same normal $(1,1,1)$, but $\frac52 \neq 1$:
strictly parallel planes. The line: direction
$(1,1,0) \cdot (1,1,1) = 2 \neq 0$: not parallel to the plane,
so it pierces it at one point.

\textbf{5.} \omterm{prop:g10:coordgeom:midpoint}{Midpoints}: $\left(1, \frac12, 0\right)$,
$\left(\frac12, 1, 0\right)$, $\left(0, 1, \frac12\right)$,
$\left(0, \frac12, 1\right)$, $\left(\frac12, 0, 1\right)$,
$\left(1, 0, \frac12\right)$: each has coordinate sum
$\frac32$.

\textbf{6.} $P$'s normal is $(1,1,1)$, the direction of
$(AG)$: perpendicular. The center
$\left(\frac12,\frac12,\frac12\right)$ has coordinate sum
$\frac32$: on $P$.

\textbf{7.} Consecutive \omterm{prop:g10:coordgeom:midpoint}{midpoints} differ by \omterm{def:g11:vect:vector}{vectors} like
$\left(-\frac12, \frac12, 0\right)$: length
$\frac{\sqrt2}{2}$ each time --- six equal sides.

\textbf{8.} At $\left(\frac12, 1, 0\right)$: \omterm{def:g11:vect:vector}{edge-vectors}
$\left(\frac12, -\frac12, 0\right)$ and
$\left(-\frac12, 0, \frac12\right)$: dot product $-\frac14$,
\omterm{def:g11:scal:dot}{norms} $\frac{\sqrt2}{2}$: $\cos = -\frac12$: angle
$120^\circ$. Equal sides, equal $120^\circ$ angles: a regular
hexagon.

\textbf{9.} Perimeter $6 \times \frac{\sqrt2}{2} = 3\sqrt2
\approx 4.24$. Area $6 \times \frac{\sqrt3}{4} \times
\frac12 = \frac{3\sqrt3}{4} \approx 1.30$: larger than a face
(area $1$) --- a slice bigger than any side of the box; only
the diagonal rectangle ($\sqrt2 \approx 1.41$) beats it.

\textbf{10.} $(t,t,t)$ meets $P$ at $t = \frac12$: the center
--- which is the \omterm{prop:g10:coordgeom:midpoint}{midpoint} of $[AG]$ ($A$ at $t = 0$, $G$ at
$t = 1$).

\textbf{11.} For $c = \frac12$: the plane cuts the three edges
at $A$'s corner, at $\left(\frac12,0,0\right)$,
$\left(0,\frac12,0\right)$, $\left(0,0,\frac12\right)$: an
equilateral triangle of side $\frac{\sqrt2}{2}$. As $c$ grows,
the triangle grows, its corners get truncated into a hexagon
(regular exactly at $c = \frac32$), then the picture shrinks
symmetrically back to a triangle at $G$'s corner: the
crystal-cut morphing.

\textbf{12.} Each side of the section is the \omterm{def:g10:numbers:interunion}{intersection} of
the cutting plane with one \emph{face} of the cube, and a
plane meets a face (a flat polygon) in at most one segment:
at most $6$ sides. Seven is impossible; three through six all
occur (questions 8 and 11 show two of them).

\textbf{13.} Base $ABD$: right triangle of area $\frac12$;
height $AE = 1$: volume
$\frac13 \times \frac12 \times 1 = \frac16$.

\textbf{14.} Every pair among $B, D, E, G$ differs in exactly
two \omterm{def:g10:coordgeom:system}{coordinates} by $1$: all six distances equal $\sqrt2$ ---
regular. The cube splits into $BDEG$ plus four corner
tetrahedra congruent to $ABDE$: volume
$1 - 4 \times \frac16 = \frac13$.

\textbf{15.} Plane $(BDE)$: $x + y + z = 1$; center
$\left(\frac12,\frac12,\frac12\right)$: distance
$\frac{\abs{\frac32 - 1}}{\sqrt3} = \frac{1}{2\sqrt3} =
\frac{\sqrt3}{6}$.

\textbf{16.} $\vect{AG} = (1,1,1)$,
$\vect{BH} = (-1,1,1)$: $\cos\theta = \frac{-1 + 1 + 1}{3} =
\frac13$: $\theta \approx 70.5^\circ$, supplement
$\approx 109.5^\circ$. Four electron pairs repel one another
and spread as far apart as possible around the carbon: the
optimum is the regular tetrahedron's corners, whose
center-to-vertex directions meet at exactly this angle ---
diamond is question 14 crystallized.

\textbf{17.} $(1, t, t)$ on $P$: $1 + 2t = \frac32$:
$t = \frac14$: the point $\left(1, \frac14, \frac14\right)$.

\textbf{18.} $(AB)$: points $(t, 0, 0)$; $(EG)$: points
$(s, s, 1)$. Meeting would need $0 = 1$ in the third
coordinate: never. Parallel would need $(1,0,0)$ and
$(1,1,0)$ \omterm{def:g12:space:collinear}{collinear}: no. Neither meeting nor parallel: skew
--- two corridors on different floors.

\textbf{19.} $B + t(1,1,1) = (1 + t, t, t)$ on $P$:
$1 + 3t = \frac32$: $t = \frac16$: projection
$\left(\frac76, \frac16, \frac16\right)$, whose coordinate sum
is $\frac32$: on $P$, as required.

\textbf{20.} Parametric lines pierce (questions 1, 10, 17);
\omterm{def:g12:space:normal}{normal vectors} measure angles and distances (questions 2, 8,
15, 16); plane \omterm{def:g10:algebra:equation}{equations} carve sections (the triangle-to-
hexagon morphing). And the playground repaid the visit: a
regular hexagon in a box of squares, a regular tetrahedron in
the corners, the angle of diamond, and two lines that ignore
each other --- geometry only space can offer.
\end{solution}
